% =============================================================================
% 050 / 068
% Status: raw
% =============================================================================
% Notes:
%
% Source question:
% 這樣理解可以嗎?
% 文讀就像是 源自 該時代官話 的外來語，
% 白讀就是 該時代既有的詞語。(白讀不等於固有詞)
% 
% 而在東亞大陸上，各語言屢次遭受各朝代官話滲透，
% 前一個朝代的文讀，會成為後一個時代的白讀。
% 可想而知，一個漢字的讀音會受各朝代影響，層層累加。
% 
% 這就是梳理構詞的困難點嗎？
% =============================================================================

\chapter[這樣理解可以嗎? 文讀就像是 源自 該時代官話 的外來語， 白讀就是 該時代既有的詞語。(白讀不等於固有詞) 而在東亞大陸上，各語言屢次遭受各朝代官話滲透， 前一個朝代的文讀，會成為後一個時代的白讀。 可想而知，一個漢字的讀音會受各朝代影響，層層累加。 這就是梳理構詞的困難點嗎？]{這樣理解可以嗎? \\[0.35em] 文讀就像是 源自 該時代官話 的外來語， \\[0.35em] 白讀就是 該時代既有的詞語。(白讀不等於固有詞) \\[0.35em] 而在東亞大陸上，各語言屢次遭受各朝代官話滲透， \\[0.35em] 前一個朝代的文讀，會成為後一個時代的白讀。 \\[0.35em] 可想而知，一個漢字的讀音會受各朝代影響，層層累加。 \\[0.35em] 這就是梳理構詞的困難點嗎？}
\qaanswer{
基本上就是這麼一回事
}

